% ====================================================================
% abstract.tex (dissertation example)
% ====================================================================

\phantomsection
\section*{Abstract}
\addcontentsline{toc}{section}{Abstract}
\label{sec:abstract}

Low-cost microcontrollers usually include a 12-bit analogue-to-digital
converter, which is not enough for sensors whose output changes by only a
small fraction of the input range. Oversampling and averaging is a well-known
way to obtain more resolution from the same converter, at the cost of a lower
output rate, but it only works if the conversion errors differ from one sample
to the next, and the amount of noise needed for that is often stated loosely.

This dissertation examines the technique by simulation. An ideal 12-bit
converter was modelled with Gaussian noise added at its input, and the
effective resolution after averaging was measured for oversampling ratios
from 1 to 256, with both a sine wave and static inputs. With
\qty{0.5}{LSB} rms of noise the simulated resolution followed the theoretical
gain of one bit for every factor of four in oversampling ratio to within
0.03~bits, reaching 15.0 effective bits at a ratio of 256. With no noise,
averaging a static input gave no improvement at all, and with
\qty{0.25}{LSB} rms the improvement stopped at around 14 bits.

The results show that about half an LSB of noise is enough for the method to
work as predicted up to a ratio of 256, and that adding more than this costs
resolution. The model does not include converter non-linearity, which
averaging cannot correct, so the figures are an upper limit for a real device.
